\documentclass[11pt]{article}
\usepackage[margin=1in]{geometry}
\usepackage[T1]{fontenc}
\usepackage[utf8]{inputenc}
\usepackage{amsmath,amssymb,amsthm}
\usepackage{graphicx}
\usepackage{xcolor}
\usepackage{microtype}
\usepackage[colorlinks=true,linkcolor=blue!55!black,citecolor=blue!55!black,urlcolor=blue!55!black]{hyperref}

\newtheorem{theorem}{Theorem}
\newtheorem{lemma}{Lemma}
\newtheorem{corollary}{Corollary}
\newcommand{\Id}{\mathrm{Id}}
\newcommand{\R}{\mathbb R}

\title{A noncommuting singular-commutator class satisfying every matrix rearrangement inequality}
\author{Candidate counterexample packet for arXiv:1904.05239}
\date{11 August 2026}

\begin{document}
\maketitle

\begin{center}
\fbox{\parbox{0.92\textwidth}{\textbf{Status: candidate full counterexample to necessity, likely valid, subject to human review.}
There are noncommuting positive definite $3\times3$ matrices $A,B$ for which
$\ker(AB-BA)\ne\{0\}$, but every rearrangement inequality in the source
holds, globally and also for $(\Id+\varepsilon A,\Id+\varepsilon B)$ for every
$\varepsilon>0$.  The same argument gives a simple sufficient condition:
a common orthogonal reducing decomposition into blocks of dimension at most
two.}}
\end{center}

\section{Source question}

Alaifari, Cheng, Pierce, and Steinerberger prove that for positive
semidefinite $2\times2$ matrices $C,D$, every word satisfies
\begin{equation}\label{eq:source2x2}
 \|C^{m_1}D^{n_1}\cdots C^{m_s}D^{n_s}\|
 \le \|C^mD^n\|,
 \qquad m=\sum_i m_i,\quad n=\sum_i n_i.
\end{equation}
They also prove a near-identity theorem in arbitrary dimension under the
condition $\ker(AB-BA)=\{0\}$ \cite[Theorems 1 and 2]{source}.  Immediately
after Theorem 2, on page 3, they write:
\begin{quote}
``We do not know whether the condition $\ker(AB-BA)=\emptyset$ is necessary
but are inclined to think that it may not be.''
\end{quote}
Here and below the printed $\emptyset$ is read as the zero subspace.  They
also ask for simple sufficient conditions other than commutativity.  The
source page is reproduced in Appendix~\ref{app:evidence}.

The literal commuting example already shows that the condition is not
necessary, and the source explicitly notes this trivial case.  The result
below removes that loophole: its matrices do not commute, their commutator
has a nonzero kernel, and the conclusion holds for every word and every
positive perturbation parameter.

\section*{Idea of the proof}

The source's $2\times2$ theorem can be assembled blockwise.  Products and
powers respect a common orthogonal direct-sum decomposition, and the spectral
norm of a direct sum is the maximum of the block norms.  Hence every pair
whose common reducing blocks have size at most two satisfies all word
inequalities.  Taking one genuinely noncommuting $2\times2$ block and adjoining
a scalar block makes the full pair noncommuting while forcing a kernel in the
commutator.  Choosing the scalar value larger than both block norms even puts
the top eigenspace of every $mA+nB$ inside that kernel.

\section{Block stability}

For nonnegative integers $m_i,n_i$ as in the source, write
\[
 W(C,D)=C^{m_1}D^{n_1}\cdots C^{m_s}D^{n_s}.
\]

\begin{lemma}[direct-sum closure]\label{lem:sum}
Suppose
\[
 C=\bigoplus_{j=1}^r C_j,\qquad D=\bigoplus_{j=1}^r D_j
\]
are positive semidefinite matrices on a common orthogonal decomposition.  If
\[
 \|W(C_j,D_j)\|\le \|C_j^mD_j^n\|
\]
for every $j$, then
\[
 \|W(C,D)\|\le \|C^mD^n\|.
\]
\end{lemma}

\begin{proof}
Functional calculus and multiplication are blockwise, so
\[
 W(C,D)=\bigoplus_j W(C_j,D_j),\qquad
 C^mD^n=\bigoplus_j C_j^mD_j^n.
\]
The spectral norm of an orthogonal direct sum is the maximum of the block
norms.  Therefore
\[
 \|W(C,D)\|=\max_j\|W(C_j,D_j)\|
 \le \max_j\|C_j^mD_j^n\|=\|C^mD^n\|.
\]
\end{proof}

\begin{corollary}[a simple noncommutative sufficient condition]\label{cor:block}
If two real symmetric positive semidefinite matrices have a common
orthogonal reducing decomposition into blocks of dimension at most two, then
they satisfy every rearrangement inequality \eqref{eq:source2x2}.
\end{corollary}

\begin{proof}
On a one-dimensional block the two matrices commute, so equality holds.  On
a two-dimensional block, inequality \eqref{eq:source2x2} is Theorem 1 of the
source.  Apply Lemma~\ref{lem:sum}.
\end{proof}

\section{Explicit counterexample to necessity}

\begin{theorem}\label{thm:example}
Let
\[
 A=\begin{pmatrix}2&0&0\\0&1&0\\0&0&4\end{pmatrix},\qquad
 B=\begin{pmatrix}2&1&0\\1&2&0\\0&0&4\end{pmatrix}.
\]
Then:
\begin{enumerate}
\item $A$ and $B$ are positive definite and do not commute;
\item $\ker(AB-BA)=\operatorname{span}\{e_3\}$;
\item for every admissible word,
\[
 \|W(A,B)\|\le \|A^mB^n\|;
\]
\item for every $\varepsilon>0$ and every admissible word,
\[
 \|W(\Id+\varepsilon A,\Id+\varepsilon B)\|
 \le
 \|(\Id+\varepsilon A)^m(\Id+\varepsilon B)^n\|.
\]
Thus the commutator-kernel condition in the source's near-identity theorem is
not necessary, even among noncommuting positive definite $3\times3$ pairs.
\end{enumerate}
\end{theorem}

\begin{proof}
Set
\[
 A_0=\begin{pmatrix}2&0\\0&1\end{pmatrix},\qquad
 B_0=\begin{pmatrix}2&1\\1&2\end{pmatrix}.
\]
The eigenvalues of $A_0$ are $2,1$, and those of $B_0$ are $3,1$; hence
$A=A_0\oplus[4]$ and $B=B_0\oplus[4]$ are positive definite.  Direct
multiplication gives
\[
 AB-BA=
 \begin{pmatrix}0&1&0\\-1&0&0\\0&0&0\end{pmatrix}.
\]
Thus $A,B$ do not commute and the kernel of their commutator is exactly
$\operatorname{span}\{e_3\}$.

The pair has a common orthogonal reducing decomposition $\R^3=\R^2\oplus
\R e_3$ with block dimensions two and one.  Part (3) follows from
Corollary~\ref{cor:block}.  For part (4), observe that
\[
 \Id+\varepsilon A=(I_2+\varepsilon A_0)\oplus[1+4\varepsilon],\qquad
 \Id+\varepsilon B=(I_2+\varepsilon B_0)\oplus[1+4\varepsilon].
\]
Both two-dimensional blocks are positive definite for $\varepsilon>0$, so
the same corollary applies.  Equivalently, on the scalar block the two sides
are equal to $(1+4\varepsilon)^{m+n}$, while the source's $2\times2$ theorem
controls the other block; taking the maximum preserves the inequality.
\end{proof}

There is also a robustness check against a natural weakening of the printed
hypothesis.  If $V_1$ denotes the top eigenspace of $mA+nB$ and $m,n\ge1$,
then on $e_3$ the eigenvalue is $4(m+n)$, whereas
\[
 \|mA_0+nB_0\|\le 2m+3n<4(m+n).
\]
Consequently $V_1=\operatorname{span}\{e_3\}\subset\ker(AB-BA)$.  Thus the
example also violates the weaker exclusion
$V_1\cap\ker(AB-BA)=\{0\}$ while retaining all rearrangement inequalities.

\section{Verification, scope, and novelty}

The formal proof uses only block multiplication, the maximum formula for the
spectral norm of an orthogonal direct sum, and Theorem 1 of the source.  The
matrix arithmetic and random finite word tests are independently checked by
\texttt{code/verify\_example.py}; those computations are sanity checks, not
part of the proof.

This answers the source's necessity question in the negative under its
nontrivial, noncommuting interpretation and supplies the adjacent requested
simple sufficient condition.  It does not address the source's broad
measure-theoretic question about generic admissibility, nor does it classify
all pairs satisfying every word inequality.

The bounded novelty check on 11 August 2026 searched the run registry and its
solution, attempt, and proof-gap indexes for the arXiv id, title, commutator
kernel, and block-rearrangement phrases.  Exact and close-phrase primary-source
web searches found the source preprint, its published record, and author
publication pages, but no later paper resolving this necessity question.
Because the argument is elementary, this is not an exhaustive priority claim.

\textbf{Human-review recommendation.} Verify that the authors intended
``necessary'' to exclude not only commuting pairs but also reducible
noncommuting pairs.  Under that natural reading, Theorem~\ref{thm:example} is
decisive; if irreducibility was intended but unstated, the irreducible variant
remains open.

\appendix
\section{Source-page evidence}\label{app:evidence}

\begin{figure}[p]
\centering
\includegraphics[width=0.94\textwidth]{figures/open_problem_crop.png}
\caption{Page 3 of arXiv:1904.05239v2.  Theorem 1 gives the $2\times2$ input;
the text following Theorem 2 asks whether the commutator-kernel condition is
necessary and asks for other simple sufficient conditions.}
\end{figure}

\begin{thebibliography}{9}
\bibitem{source}
R. Alaifari, X. Cheng, L. B. Pierce, and S. Steinerberger,
\emph{On Matrix Rearrangement Inequalities},
Proc. Amer. Math. Soc. 148 (2020), 1835--1848;
arXiv:1904.05239v2 [math.FA].
\end{thebibliography}

\end{document}
